\scanpage{366}
\mspage{90}
nullement clair que l'intersection de deux localiseurs fond. satisfaisant à $\mathrm{Loc}(7\mathrm{F})$ y satisfasse.

$\mathrm{Loc}(7\mathrm{D}_{W_{\unsure{\omega|w}}})$ est exactement ce qu'il faut pour assurer que la factorisation canonique d'une flèche $f:X\unsure{-|\to}Y$ de $\mathrm{Cat}$ en $f=pi$, avec $i\in W_\omega$, $p\in\mathrm{Parf}_{W_\omega}$, est aussi une factorisation \deleted{\unsure{de}} \unsure{du} type stipulé par l'axiome de factorisation $\mathrm{Loc}(7\mathrm{F})$ : \unsure{on a|car.} (bien sûr) $i\in W$, mais aussi $p\in\mathrm{Fib}_W$ (car $\mathrm{D}_{W_\omega}$ équivaut à sa forme forte $\mathrm{Parf}_{W_\omega}\subset\mathrm{Fib}_W$). Cela signifie aussi que la factorisation canonique de $f$ peut être utilisée pour ``calculer'' ses fibres homotopiques, non seulement en les points de $Y$, mais aussi \unsure{p.r.} à un morphisme de chgt. de base $Y'\unsure{\longrightarrow|\Longrightarrow}Y$. \deleted{\unsure{?}}

En l'absence de cet axiome,
% archive sheet number: 365
